% !TeX root = ../notes.tex

\chapter*{About this course}
\addcontentsline{toc}{chapter}{About this course}

\section*{What this course will achieve}
This course connects physical organic chemistry, cheminformatics, data analysis, reaction optimisation and laboratory automation. Its central question is how chemists can translate molecules and experiments into structured data, extract chemical understanding from those data, and use that understanding to choose and perform better experiments.

\section*{What is digital chemistry?}
In this course, \emph{digital chemistry} means the use of machine-readable molecular and experimental information to support chemical understanding and experimental decision-making. It includes molecular representations, chemical-space analysis, quantitative physical organic methods, machine learning, optimisation, automated experimentation and closed-loop laboratories. It is not treated as a replacement for chemical reasoning: the objective is to make chemical reasoning quantitative, reproducible and scalable.

\begin{warningbox}
\textbf{Scope:} digital chemistry is broader than artificial intelligence, and automation alone does not make an experiment autonomous.
\end{warningbox}

\section*{Learning outcomes}
By the end of the course, students should be able to:
\begin{enumerate}
  \item compare common molecular representations;
  \item explain how descriptors define different chemical spaces;
  \item interpret VTNA and RPKA plots mechanistically;
  \item distinguish regression, classification, and optimisation problems;
  \item assess model validation, and chemical usefulness;
  \item compare OVAT, design of experiments and Bayesian optimisation;
  \item explain how to generate structured chemical data;
  \item design a conceptual closed-loop workflow that balances chemical performance, sustainability and experimental constraints.
\end{enumerate}

\section*{Case-study threads}
\begin{itemize}
  \item \textbf{Physical organic chemistry:} asymmetric catalysis, non-linear effects, kinetics, RPKA and VTNA.
  \item \textbf{Green chemistry:} solvent selection, process mass intensity, multi-objective optimisation and responsible automation.
  \item \textbf{Cheminformatics:} molecular representations, fingerprints, similarity, descriptors and molecular libraries.
\end{itemize}

\begin{insightbox}
\textbf{Recurring question:} what information is needed, how should it be represented, and what chemical decision will it support?
\end{insightbox}

\begin{conceptbox}
  Concept box
\end{conceptbox}

\begin{warningbox}
  Warning box
\end{warningbox}

\begin{methodbox}
  Method box
\end{methodbox}

\begin{resultbox}
  Result box
\end{resultbox}

\begin{notebox}
  Note box
\end{notebox}